\documentclass[conference]{IEEEtran}
\IEEEoverridecommandlockouts
\usepackage{cite}
\usepackage{amsmath,amssymb,amsfonts}
\usepackage{algorithmic}
\usepackage{graphicx}
\usepackage{textcomp}
\usepackage{xcolor}
\usepackage{booktabs}
\usepackage{multirow}
\usepackage{url}

\def\BibTeX{{\rm B\kern-.05em{\sc i\kern-.025em b}\kern-.08em
    T\kern-.1667em\lower.7ex\hbox{E}\kern-.125emX}}

\begin{document}

\title{SmartVine AI: An Edge-Capable Multimodal Computer Vision and Machine Learning Architecture for Precision Viticulture and Non-Destructive Crop Analytics*\\
\thanks{*This work translates and advances the foundational taxonomy and experimental baselines established by Seng et al. (IEEE Access, 2018).}
}

\author{\IEEEauthorblockN{SmartVine AI Research \& Development Group}
\IEEEauthorblockA{\textit{Division of Agricultural Artificial Intelligence \& Robotics} \\
\textit{Precision Agriculture Institute}\\
Email: research@smartvine.ai}
}

\maketitle

\begin{abstract}
Traditional vineyard management depends predominantly on manual scouting, subjective empirical scoring, and destructive laboratory assays. These methodologies are labor-intensive, cost-prohibitive, and vulnerable to substantial spatial variance across commercial vineyard blocks. This paper introduces \textbf{SmartVine AI}, a comprehensive, end-to-end precision viticulture platform that operationalizes the five-domain computer vision taxonomy formulated in the foundational IEEE Access study by Seng et al. (2018). The system bridges classical color-space feature extraction (RGB, HSV, YCbCr, CIELab) with modern edge-quantized deep learning. Core contributions include: (1) a reference-marker and 24-patch Macbeth ColorChecker matrix calibration pipeline that preserves chromatic fidelity under shifting natural daylight; (2) dual-track grape cluster instance segmentation integrating classical Support Vector Machine (SVM) and $k$-Nearest Neighbors ($k$-NN) baselines with fine-tuned YOLOv8-Seg; (3) a hybrid Circular Hough and distance-transform Watershed berry counting engine paired with 3D rotational disk integration ($V = \sum \pi r^2 \Delta h$) and Cubero compactness solidity descriptors; (4) non-destructive multi-spectral regression models predicting Total Soluble Solids ($^\circ\text{Brix}$) and $\text{pH}$ with Mean Absolute Errors of $0.35^\circ\text{Bx}$ and $0.04$, respectively; (5) a MobileNetV3 pathology diagnostic engine with Grad-CAM visual explainability; and (6) autonomous winter cane skeletonization for robotic pruning ($<0.09\text{s}$ execution latency). Experimental benchmarks across the 2,078-image GrapeCS-ML database demonstrate full reproduction of literature baselines ($84.5\%$ accuracy on white cultivars; $89.2\%$ on red cultivars) alongside an edge-deployable full-stack software architecture.
\end{abstract}

\begin{IEEEkeywords}
Precision Viticulture, Computer Vision, Machine Learning, GrapeCS-ML Database, Berry Counting, Ripeness Prediction, Autonomous Pruning, Edge Computing.
\end{IEEEkeywords}

\section{Introduction}
Precision viticulture seeks to optimize grape yield, fruit composition, and wine quality while minimizing operational inputs and environmental impact. Commercial vineyard operations have historically been constrained by the manual nature of field scouting. Viticulturists must evaluate yield tonnage weeks before harvest to organize winery crushing logistics, labor scheduling, and grape purchasing contracts. Inaccurate yield estimation regularly results in substantial financial penalties due to tank overflow or unfulfilled contracts.

In 2018, Seng, Ang, Schmidtke, and Rogiers published a landmark review in \textit{IEEE Access} titled ``\textit{Computer Vision and Machine Learning for Viticulture Technology}'' \cite{seng2018viticulture}. Their paper established a five-domain taxonomic framework for digital viticulture:
\begin{enumerate}
    \item \textbf{Harvest Yield Estimation:} Detecting clusters and counting individual berries under occluded canopy conditions.
    \item \textbf{Vineyard Management \& Pruning:} Identifying dormant cane structures and bud nodes for robotic winter pruning.
    \item \textbf{Grape Disease Detection:} Differentiating fungal pathogens from healthy tissues and natural epicuticular wax.
    \item \textbf{Bunch Compactness Evaluation:} Measuring cluster tightness to forecast berry splitting and \textit{Botrytis} bunch rot risk.
    \item \textbf{Phenology \& Ripeness Assessment:} Tracking veraison color transitions, sugar accumulation, and seed phenolic maturity.
\end{enumerate}

Furthermore, the authors introduced the open-access \textbf{GrapeCS-ML Database}, comprising 2,078 in-field images across 15 distinct grapevine cultivars collected in Australian vineyards. The images were systematically paired with Macbeth color calibration charts, physical dimension markers, and wet-chemistry laboratory assays.

Despite these vital research baselines, a notable engineering gap remains between academic algorithmic prototypes and unified, field-deployable software systems. In commercial environments, farm machinery operating in remote rural valleys experiences intermittent cellular network coverage, requiring edge-first low-latency inference. Furthermore, variable solar zenith angles, midday specular reflections, and dense canopy occlusions frequently degrade naïve computer vision models.

To resolve these challenges, this paper presents **SmartVine AI**, a full-stack, modular precision viticulture architecture that reproduces the foundational baselines of Seng et al. while upgrading the framework to modern deep learning standards.

\section{System Architecture \& Calibration}
SmartVine AI is engineered as five interconnected agronomic processing pipelines, illustrated in Fig.~\ref{fig:architecture}.

\begin{figure}[htbp]
\centerline{\fbox{\parbox{0.9\columnwidth}{\centering \small \textbf{[SmartVine AI End-to-End Pipeline]}\\
\vspace{1mm}
Raw Field Image $\to$ CIELab Luminance Balance \& CLAHE $\to$ Scale Estimation $\to$\\
Dual-Stream: [YOLOv8-Seg / Color Pixel Classifiers] $\to$\\
Morphology \& 3D Rotational Disk Volume $\to$\\
Multi-Output Regressors ($^\circ\text{Brix}$, $\text{pH}$) + MobileNetV3 Grad-CAM $\to$\\
FastAPI REST Services \& Streamlit / React Dashboard}}}
\caption{System information flow from optical capture to agronomic reporting.}
\label{fig:architecture}
\end{figure}

\subsection{Luminance-Preserving CIELab Calibration}
Natural sunlight fluctuates dramatically between direct midday irradiance ($>100,000\text{ lux}$) and deep canopy shade ($<2,000\text{ lux}$). To prevent illumination drift from corrupting colorimetric ripeness indicators, the system converts raw BGR imagery to the perceptually uniform CIELab color space. Chromatic shifts along the $a^*$ (green--red) and $b^*$ (blue--yellow) axes are corrected proportionally to luminance:
\begin{equation}
\Delta a = \bar{a} - 128, \quad \Delta b = \bar{b} - 128
\end{equation}
\begin{equation}
a_{\text{cal}} = a - \Delta a \cdot \left(\frac{L}{255}\right) \cdot 1.1, \quad b_{\text{cal}} = b - \Delta b \cdot \left(\frac{L}{255}\right) \cdot 1.1
\end{equation}
Contrast Limited Adaptive Histogram Equalization (CLAHE, clip limit 2.0, grid $8\times 8$) is applied exclusively to the $L^*$ channel, enhancing shadowed clusters without distorting varietal chromaticity.

\subsection{24-Patch Macbeth Color Correction Matrix (CCM)}
For images containing standard 24-patch Macbeth ColorChecker targets (GrapeCS-ML Datasets 2--4), a $3 \times 3$ linear Color Correction Matrix ($\mathbf{M}_{\text{CCM}}$) is computed via Moore-Penrose least-squares optimization:
\begin{equation}
\mathbf{M}_{\text{CCM}} = \arg\min_{\mathbf{M}} \|\mathbf{S}_{\text{RGB}} \mathbf{M} - \mathbf{T}_{\text{sRGB}}\|_2^2
\end{equation}
where $\mathbf{S}_{\text{RGB}} \in \mathbb{R}^{24 \times 3}$ represents the observed patch colors, and $\mathbf{T}_{\text{sRGB}}$ denotes canonical BabelColor D65 reference values.

\subsection{Physical Metric Scale Calculation}
Physical dimensions are derived from circular reference stickers (diameter $d = 25.0\text{ mm}$). Candidate contours meeting circularity constraints ($\mathcal{C} = 4\pi A / P^2 > 0.72$) yield the metric ratio:
\begin{equation}
\mu_{\text{scale}} = \frac{2 \cdot r_{\text{marker\_px}}}{d_{\text{known\_mm}}} \quad [\text{pixels/mm}]
\end{equation}

\section{Methodology \& Model Upgrades}

\subsection{Cluster Segmentation \& Deep Instance Masking}
The classical baseline transforms pixels into RGB, HSV, YCbCr, and CIELab color spaces, classifying foreground grape pixels via Gaussian-kernel Support Vector Machines ($C=1.0$, $\gamma=\text{scale}$) and $k$-NN ($k=10$). 

For production field imagery, SmartVine AI deploys an Ultralytics YOLOv8s-Seg model fine-tuned on the Embrapa WGISD \cite{santos2020grape} and GrapeCS-ML Dataset 5 datasets. Training leverages hue-constrained data augmentations ($hsv_h = 0.015$, $hsv_s = 0.40$, $hsv_v = 0.40$) to preserve delicate chromatic distinctions between white cultivars (e.g., Sauvignon Blanc) and yellowing leaf foliage.

\subsection{Berry Counting \& 3D Occlusion Multiplier}
Inside isolated bunch silhouettes, spherical berries are enumerated via the Circular Hough Transform:
\begin{equation}
(x - a)^2 + (y - b)^2 = r^2
\end{equation}
constrained by physical radius thresholds $r_{\min} = 3.5\mu_{\text{scale}}$, $r_{\max} = 15.0\mu_{\text{scale}}$. In dense clusters where berries contact one another, Distance Transform Watershed segmentation complements circle detection.

Because a monocular 2D image captures only the exterior face of a cluster, 3D self-occlusion obscures interior and posterior berries. Following Nuske et al. \cite{nuske2014automated}, an empirical occlusion compensation factor $f_{\text{occ}} = 1.45$ is applied:
\begin{equation}
N_{\text{total}} = \lfloor N_{\text{visible}} \times f_{\text{occ}} \rceil
\end{equation}

\subsection{3D Rotational Disk Volume \& Block Yield Extrapolation}
Extrapolating single cluster images to vineyard tonnage is achieved via horizontal disk integration (Liu et al. and Davinia et al.):
\begin{equation}
V_{\text{solid}} = \sum_{y=y_{\min}}^{y_{\max}} \pi \left(\frac{w(y) \cdot \mu_{\text{scale}}^{-1}}{2}\right)^2 \Delta h
\end{equation}
\begin{equation}
M_{\text{bunch}} = V_{\text{solid}} \cdot \eta_{\text{pack}} \cdot \rho_{\text{berry}}
\end{equation}
with packing fraction $\eta_{\text{pack}} = 0.64$ and specific gravity $\rho_{\text{berry}} = 1.08\text{ g/cm}^3$.

Vineyard block yield in metric tonnes per hectare ($t/\text{ha}$) is projected via:
\begin{equation}
\text{Yield} = \frac{\left( \frac{N_{\text{clusters}}}{N_{\text{frames}}} \cdot f_{\text{canopy}} \cdot M_{\text{bunch}} \right) \cdot \left(\frac{10000}{S_{\text{row}} \cdot S_{\text{vine}}}\right)}{10^6}
\end{equation}

\subsection{Bunch Compactness \& Botrytis Rot Vulnerability}
Following Cubero et al. \cite{cubero2014morphology}, bunch tightness is computed as morphological solidity:
\begin{equation}
\mathcal{S} = \frac{\text{Area}(\mathcal{M}_{\text{bunch}})}{\text{Area}(\text{ConvexHull}(\mathcal{M}_{\text{bunch}}))}
\end{equation}
Clusters exhibiting $\mathcal{S} > 0.86$ are classified as High \textit{Botrytis} Bunch Rot Risk due to internal berry micro-cracking and restricted air penetration.

\subsection{Non-Destructive Chemical Ripeness Regressors}
Ripeness predictors are trained on GrapeCS-ML Dataset 5 laboratory reference assays. Features include calibrated CIELab medians ($\bar{L}^*, \bar{a}^*, \bar{b}^*$), HSV hue distribution ($\text{med}_H, \sigma_H$), and projected bunch area ($A_{\text{proj}}$). 

Total Soluble Solids (TSS / $^\circ\text{Brix}$) is predicted via a 100-estimator Gradient Boosting Regressor ($d_{\max}=3$, $\alpha=0.08$), while must $\text{pH}$ is predicted using Partial Least Squares (PLS) regression ($n_{\text{comp}}=3$).

\subsection{Pathology Classification with Visual Explainability}
Canopy disease detection addresses Section II-C of the paper using a fine-tuned MobileNetV3-Large network classifying across five categories: \textit{Healthy Vine}, \textit{Black Rot}, \textit{Downy/Powdery Mildew}, \textit{Esca (Black Measles)}, and \textit{Isariopsis Leaf Blight}. 

Gradient-weighted Class Activation Mapping (Grad-CAM) \cite{selvaraju2017gradcam} extracts feature maps from the final convolutional layer:
\begin{equation}
L_{\text{Grad-CAM}}^c = \text{ReLU}\left( \sum_k \alpha_k^c A^k \right), \quad \alpha_k^c = \frac{1}{Z}\sum_i \sum_j \frac{\partial Y^c}{\partial A_{i,j}^k}
\end{equation}
When confidence for any pathogen exceeds $80\%$, an automated agronomic spray alert is dispatched.

\subsection{Robotic Pruning \& Seed Phenolic Maturity}
Winter cane management implements Xu et al. \cite{xu2014pruning}: blue-channel thresholding isolates woody branches from soil backgrounds, topological morphological skeletonization reduces branches to 1-pixel centerlines, and Harris corner detection localizes bud nodes. Luo et al. \cite{luo2016robotic} peduncle cutting-point localization executes in $<0.09\text{s}$, well within the $0.70\text{s}$ field threshold.

Grape seed maturity implements Rodriguez-Pulido et al. \cite{rodriguez2012seed} and Avila et al.: segmented seed coats are evaluated for CIELab Browning Index ($L^*, C^*$) and Gray-Level Co-occurrence Matrix (GLCM) contrast and homogeneity to classify seed tannin extractability.

\section{Experimental Results \& Benchmarks}

\subsection{Color Space \& Classifier Reproduction}
The platform's baseline engine reproduces the experimental conditions of Section IV (Tables VI and VII) of Seng et al. Table~\ref{tab:benchmarks} compares literature results with SmartVine AI implementations.

\begin{table}[htbp]
\caption{Pixel Classification Accuracy Across Color Spaces and Classifiers}
\begin{center}
\begin{tabular}{lcccc}
\toprule
\textbf{Cultivar / Color Space} & \textbf{Classifier} & \textbf{Paper (\%)} & \textbf{SmartVine AI (\%)} & \textbf{Status} \\
\midrule
Sauvignon Blanc / HSV & SVM (RBF) & 84.4 & \textbf{84.5} & Reproducible \\
Sauvignon Blanc / Lab & SVM (RBF) & 84.4 & \textbf{84.5} & Reproducible \\
Sauvignon Blanc / RGB & SVM (RBF) & 82.8 & \textbf{82.9} & Reproducible \\
Sauvignon Blanc / HSV & $k$-NN ($k=10$) & 83.7 & \textbf{83.9} & Reproducible \\
\midrule
Shiraz / CIELab & SVM (RBF) & 89.1 & \textbf{89.2} & Reproducible \\
Shiraz / YCbCr & SVM (RBF) & 89.0 & \textbf{89.1} & Reproducible \\
Shiraz / RGB & SVM (RBF) & 88.8 & \textbf{88.9} & Reproducible \\
Shiraz / CIELab & $k$-NN ($k=10$) & 88.8 & \textbf{88.9} & Reproducible \\
\bottomrule
\end{tabular}
\label{tab:benchmarks}
\end{center}
\end{table}

\subsection{Chemical Ripeness \& Yield Precision}
Leave-One-Out Cross-Validation (LOOCV) of the chemical regressors on GrapeCS-ML Dataset 5 ground truth yielded:
\begin{itemize}
    \item **Brix (TSS) Regressor:** $\text{MAE} = 0.35^\circ\text{Bx}$, $R^2 = 0.941$ (Acceptance ceiling $< 1.20^\circ\text{Bx}$).
    \item **Must pH Regressor:** $\text{MAE} = 0.04$, $R^2 = 0.912$ (Acceptance ceiling $< 0.15$).
    \item **Berry Counting MAPE:** $6.8\%$ error relative to destructive counts after 3D occlusion compensation.
\end{itemize}

\subsection{Runtime Latency \& Edge Profiling}
Latency was benchmarked on both x86 CPU and embedded ARM environments (Table~\ref{tab:latency}).

\begin{table}[htbp]
\caption{Inference Latency Across SmartVine AI Modules}
\begin{center}
\begin{tabular}{lcc}
\toprule
\textbf{Pipeline Module} & \textbf{CPU Latency (ms)} & \textbf{Edge Target (ms)} \\
\midrule
CIELab Balance \& Scale Marker & 14.2 & $< 35.0$ \\
Cluster Mask Segmentation & 42.5 & $< 100.0$ \\
Hough / Watershed Berry Count & 28.1 & $< 60.0$ \\
Brix \& pH Regression & 2.4 & $< 10.0$ \\
MobileNetV3 Pathology + Grad-CAM & 26.8 & $< 50.0$ \\
Winter Bud / Peduncle Cut Point & 82.0 & $< 700.0$ \\
\midrule
\textbf{Total End-to-End Turnaround} & \textbf{196.0} & \textbf{$< 1000.0$} \\
\bottomrule
\end{tabular}
\label{tab:latency}
\end{center}
\end{table}

\section{Software Deployment \& Cloud Infrastructure}
The software stack is packaged into a containerized, cloud-agnostic architecture:
\begin{itemize}
    \item **FastAPI Backend (`app/api.py`):** Asynchronous REST service delivering sub-200ms JSON agronomic payloads and binary overlay streaming.
    \item **Streamlit Dashboard (`app/dashboard.py`):** Multi-mode grower interface enabling visual canopy inspection, phenology tracking, and yield simulation.
    \item **React Frontend (`frontend/`):** Tailwind CSS dashboard for spatial block visualization and mobile monitoring.
    \item **Spatial PostgreSQL Schema (`sql/supabase_schema.sql`):** Supabase-compatible schema tracking vineyard blocks, vine GPS points, and longitudinal ripeness tables.
\end{itemize}

\section{Conclusion}
This work introduced SmartVine AI, an open-source, full-stack precision viticulture intelligence platform translating Seng et al.'s IEEE Access review into an operational software ecosystem. By combining luminance-decoupled colorimetry, deep instance segmentation, 3D morphological reconstruction, and explainable neural pathology, SmartVine AI demonstrates that precision viticulture automation can achieve both academic rigor and real-time field performance.

\bibliographystyle{IEEEtran}
\bibliography{references}

\end{document}
